\section{Powertrain Modeling and Motor Selection}

\begin{multicols}{2}
The foundation of the electric vehicle (EV) powertrain begins with accurately quantifying the opposing environmental forces to derive the requisite mechanical output of the motor. This section details the fundamental longitudinal vehicle dynamics, motor rating selection, and the implementation of advanced control algorithms required to manage the tractive power effectively.

\subsection{Longitudinal Vehicle Dynamics}

To define the continuous and peak operational limits of the electric motor, the total tractive effort ($F_{te}$) required at the tire contact patch must be explicitly calculated. This effort must overcome four distinct environmental and inertial components: Aerodynamic Drag ($F_{ad}$), Rolling Resistance ($F_{rr}$), Gradient Force ($F_{g}$), and Acceleration Force ($F_{a}$).

Assuming a standard lightweight EV geometry, the vehicle mass ($m$) is established at 300 kg, possessing a frontal area ($A$) of 1.0 m$^2$, and a target top speed ($v$) of 80 km/h (22.22 m/s).

\textbf{1. Aerodynamic Drag ($F_{ad}$):} As velocity increases, the required force to displace the fluid volume scales quadratically. Given an air density ($\rho$) of 1.225 kg/m$^3$ and a conservative drag coefficient ($C_d$) of 0.4:
\begin{equation} \label{eq:aero_drag}
\begin{split}
F_{ad} &= \frac{1}{2} \rho v^2 C_d A \\
&= 0.5 \cdot 1.225 \cdot (22.22)^2 \cdot 0.4 \cdot 1.0 \\
&\approx 121.0\,\mathrm{N}
\end{split}
\end{equation}

\textbf{2. Rolling Resistance ($F_{rr}$):} The hysteresis losses generated by the deformation of the viscoelastic tire compound against the tarmac are quantified by the coefficient of rolling resistance ($\mu_r \approx 0.015$). With gravitational acceleration ($g$) at 9.81 m/s$^2$:
\begin{equation} \label{eq:rolling_res}
\begin{split}
F_{rr} &= m \cdot g \cdot \mu_r \\
&= 300 \cdot 9.81 \cdot 0.015 \\
&= 44.1\,\mathrm{N}
\end{split}
\end{equation}

\textbf{3. Gradient/Hill Climbing Force ($F_{g}$):} To ensure operational viability on inclines, the powertrain must overcome the gravitational component acting parallel to the road surface. Assuming a design operational gradient ($\alpha$) of $5^\circ$ ($\sim$8.7\% grade):
\begin{equation} \label{eq:gradient_force}
\begin{split}
F_{g} &= m \cdot g \cdot \sin(\alpha) \\
&= 300 \cdot 9.81 \cdot \sin(5^\circ) \\
&= 256.3\,\mathrm{N}
\end{split}
\end{equation}

\textbf{4. Acceleration Force ($F_{a}$):} To achieve a target 0-80 km/h sprint in 11 seconds, an average linear acceleration ($a$) of approximately 2 m/s$^2$ is required:
\begin{equation} \label{eq:accel_force}
\begin{split}
F_{a} &= m \cdot a \\
&= 300 \cdot 2.0 \\
&= 600.0\,\mathrm{N}
\end{split}
\end{equation}

The cumulative tractive effort ($F_{te\_peak}$) under maximum acceleration on a $5^\circ$ incline is the direct summation of these vectors:
\begin{equation} \label{eq:total_tractive}
\begin{split}
F_{te\_peak} &= F_{ad} + F_{rr} + F_{g} + F_{a} \\
&= 121.0 + 44.1 + 256.3 + 600.0 \\
&= 1021.4\,\mathrm{N}
\end{split}
\end{equation}
\end{multicols}

\begin{figure}[H]
\centering
\includegraphics[width=\textwidth]{"Figures Images/figure6_foc.png"}
\caption{Field Oriented Control (FOC) Block Diagram utilized for precise regulation of IPM instantaneous torque.}
\label{fig:foc_diagram}
\end{figure}

\begin{multicols}{2}
\subsection{Motor Rating Determination}

Translating the linearly derived tractive effort into rotational motor requirements demands accounting for the mechanical drivetrain efficiency ($\eta \approx 0.90$) and the dynamic tire radius ($r = 0.25$ m).

The absolute peak power ($P_{peak\_motor}$) required at the motor shaft during maximum velocity acceleration is:
\begin{equation} \label{eq:motor_peak_pwr}
\begin{split}
P_{peak\_motor} &= \frac{F_{te\_peak} \cdot v}{\eta} \\
&= \frac{1021.4 \cdot 22.22}{0.90} \\
&\approx 25,216\,\mathrm{W} \approx 25.2\,\mathrm{kW}
\end{split}
\end{equation}

Conversely, the nominal continuous power rating ($P_{cont\_motor}$) is defined by the steady-state cruise condition on flat ground ($F_g = 0$, $F_a = 0$):
\begin{equation} \label{eq:motor_cont_pwr}
\begin{split}
P_{cont\_motor} &= \frac{(F_{ad} + F_{rr}) \cdot v}{\eta} \\
&= \frac{(121.0 + 44.1) \cdot 22.22}{0.90} \\
&\approx 4,076\,\mathrm{W} \approx 4.1\,\mathrm{kW}
\end{split}
\end{equation}

Therefore, the finalized motor specification dictates a continuous rating of 5 kW to maintain a thermal safety factor, with a peak transient overload capability of 25.2 kW.

\subsection{Motor Selection: IPM vs. SPM}

The selection between a Surface Permanent Magnet (SPM) and Interior Permanent Magnet (IPM) motor dramatically influences the torque density and high-speed field-weakening capabilities of the 72V powertrain.

While an SPM motor simplifies the rotor construction by adhering the magnets to the outer rotor circumference, it suffers from a near-homogenous magnetic reluctance path ($L_d \approx L_q$), restricting it exclusively to magnetic alignment torque. Furthermore, the fragile mechanical retention limits rotational speeds, reducing power density.

Therefore, an \textbf{Interior Permanent Magnet (IPM)} topology is explicitly selected for this powertrain. Embedding the magnets within the rotor lamination steel creates magnetic saliency ($L_q > L_d$). This asymmetric reluctance path unlocks a secondary, profound torque-producing component. The total electromagnetic torque ($T_e$) of the selected IPM motor is defined as:

\begin{equation} \label{eq:ipm_torque}
T_e = \underbrace{\frac{3}{2} p \lambda_f i_q}_{\text{Magnetic Torque}} + \underbrace{\frac{3}{2} p (L_d - L_q) i_d i_q}_{\text{Reluctance Torque}}
\end{equation}

where $p$ represents pole pairs, $\lambda_f$ is the permanent magnet flux linkage, and $i_d, i_q$ are the direct and quadrature axis stator currents. At lower base speeds, the controller biases current heavily towards $i_q$ to maximize magnetic alignment. As velocity increases, the controller manipulates $i_d$ (the negative $d$-axis current), strategically leveraging the $(L_d - L_q)$ reluctance inequality to generate supplementary torque, enabling superior overall system efficiency and high-speed performance compared directly against an SPM equivalent.

\subsection{Matching Control Method (FOC)}

Unlocking the complex reluctance torque of the IPM architecture necessitates an advanced Field Oriented Control (FOC) strategy, fundamentally superior to primitive Six-Step BLDC commutation or scalar V/Hz methods.

The FOC algorithm mathematically decouples the 3-phase AC stator currents ($i_a, i_b, i_c$) via consecutive Clarke and Park transformations into a synchronous rotating reference frame aligned with the internal rotor magnet flux. This reduces the complex AC equations into two controllable DC orthogonal components: flux-producing current ($i_d$) and torque-producing current ($i_q$).

\textbf{Maximum Torque Per Ampere (MTPA) Strategy:}
Unlike an SPM motor where MTPA strictly enforces $i_d = 0$, optimal control of the selected IPM motor requires a dynamic $i_d$ reference to maximize the combined magnetic and reluctance torque output per unit of stator current. The MTPA algorithm calculates the optimal negative $i_d$ trajectory as a function of the operational torque demand:

\begin{equation} \label{eq:ipm_mtpa}
i_d^* = \frac{\lambda_f}{2(L_q - L_d)} - \sqrt{ \frac{\lambda_f^2}{4(L_q - L_d)^2} + (i_q^*)^2 }
\end{equation}

Translation of the resultant Proportional-Integral (PI) controller reference voltages back into physical gating commands is achieved via Space Vector PWM (SVPWM). SVPWM synthesizes the rotating voltage vector by symmetrically modulating the eight fundamental switching states of the 3-phase inverter bridge. This matching topology provides a 15\% superior utilization of the 72V DC bus compared to standard sinusoidal PWM, critical for extracting the full 25.2 kW peak rating from a low-voltage battery pack.
\end{multicols}

% ============================================================
%  NEW SUBSECTION — Ground-Up EV Traction Design
%  Reference Calculations (college project EV design)
% ============================================================

\subsection{Ground-Up EV Traction Design --- Reference Calculations}

The following subsections present a complete ground-up electric vehicle traction design, computed independently from the F1 powertrain parameters above. These reference calculations validate the motor sizing methodology, battery pack architecture, DC-DC converter design, and motor control technique for a lightweight urban EV platform operating at \SI{72}{\volt} nominal.

\subsubsection{Motor Rating Calculation}

The motor rating is derived from first-principles longitudinal dynamics for a representative urban EV operating condition. The vehicle parameters are established as follows: total mass $M = \SI{1300}{\kilo\gram}$ (comprising \SI{1000}{\kilo\gram} vehicle mass and \SI{300}{\kilo\gram} passenger load), frontal width $W = \SI{1.00}{\metre}$, frontal height $H = \SI{1.50}{\metre}$, yielding a frontal area $A = \SI{1.50}{\metre\squared}$. The rolling resistance coefficient is $\mu_{rr} = 0.02$, air density $\rho = \SI{1.25}{\kilo\gram\per\metre\cubed}$, drag coefficient $C_d = 0.30$, gravitational acceleration $g = \SI{9.80}{\metre\per\second\squared}$, and the design speed is $v = \SI{40}{\kilo\metre\per\hour} = \SI{11.11}{\metre\per\second}$ on a gradient of $\theta = 10^\circ$.

The rolling resistance force is calculated as:
\begin{equation}
\begin{split}
F_{rr} &= \mu_{rr} \cdot M \cdot g \\
&= 0.02 \times 1300 \times 9.80 \\
&= 254.8\,\mathrm{N}
\end{split}
\end{equation}

The aerodynamic drag force at the design speed is:
\begin{equation}
\begin{split}
F_{aero} &= \frac{1}{2} C_d \cdot \rho \cdot A \cdot v^2 \\
&= 0.5 \times 0.30 \times 1.25 \times 1.50 \times (11.11)^2 \\
&= 34.72\,\mathrm{N}
\end{split}
\end{equation}

The gradient climbing force required to sustain motion on a $10^\circ$ incline is:
\begin{equation}
\begin{split}
F_{hill} &= M \cdot g \cdot \sin(\theta) \\
&= 1300 \times 9.80 \times \sin(10^\circ) \\
&= 1300 \times 9.80 \times 0.17365 \\
&= 2211.4\,\mathrm{N}
\end{split}
\end{equation}

The total tractive effort is the summation of all three resistive components:
\begin{equation}
\begin{split}
F_{total} &= F_{rr} + F_{aero} + F_{hill} \\
&= 254.8 + 34.72 + 2211.4 \\
&= 2501\,\mathrm{N}
\end{split}
\end{equation}

The mechanical power required at the wheels is the product of total force and velocity:
\begin{equation}
\begin{split}
P &= F_{total} \cdot v \\
&= 2501 \times 11.11 \\
&= 27{,}786\,\mathrm{W} \approx 27.8\,\mathrm{kW}
\end{split}
\end{equation}

Accounting for a drivetrain efficiency of $\eta = 0.90$, the required motor shaft power is:
\begin{equation}
\begin{split}
P_{motor} &= \frac{P}{\eta} \\
&= \frac{27.8}{0.90} \\
&= 30.9\,\mathrm{kW}
\end{split}
\end{equation}

Based on this calculation, a \SI{30}{\kilo\watt} Interior Permanent Magnet Synchronous Motor (IPMSM) rated at \SI{72}{\volt} and \SI{3000}{RPM} is selected. The IPM topology is specifically justified: the salient rotor geometry produces both magnet torque and reluctance torque simultaneously, yielding higher torque density at low speed — a critical requirement for hill-climbing on a $10^\circ$ gradient where the gradient force alone constitutes 88\% of the total tractive effort.

\subsubsection{Battery Pack Design}

The energy storage system for this \SI{72}{\volt} platform is constructed from commercial 18650-format lithium-ion cells. The pack design is presented in three parts: electrical, mechanical, and thermal.

\paragraph{Electrical Design.}
The baseline cell is a lithium-ion 18650 cylindrical cell with a nominal voltage of \SI{3.7}{\volt} and a capacity of \SI{5}{\ampere\hour}. The number of series cells required to reach the target pack voltage is:
\begin{equation}
\begin{split}
N_s &= \frac{V_{pack}}{V_{cell}} \\
&= \frac{72}{3.7} \\
&= 19.46 \rightarrow 20\,\text{cells (rounded)}
\end{split}
\end{equation}

The actual pack nominal voltage with 20 series cells is $20 \times 3.7 = \SI{74}{\volt}$, which is within 2.7\% of the \SI{72}{\volt} target and therefore acceptable.

To achieve a target pack capacity of \SI{40}{\ampere\hour}, the number of parallel cells is:
\begin{equation}
\begin{split}
N_p &= \frac{C_{pack}}{C_{cell}} \\
&= \frac{40}{5} \\
&= 8\,\text{cells in parallel}
\end{split}
\end{equation}

The total cell count is $20 \times 8 = 160$ cells in a 20S8P configuration. The total stored energy is:
\begin{equation}
\begin{split}
E_{pack} &= V_{pack} \cdot C_{pack} \\
&= 74\,\mathrm{V} \times 40\,\mathrm{Ah} \\
&= 2960\,\mathrm{Wh} \approx 3.0\,\mathrm{kWh}
\end{split}
\end{equation}

At peak \SI{30}{\kilo\watt} discharge, the pack current is:
\begin{equation}
\begin{split}
I_{pack} &= \frac{P}{V_{pack}} \\
&= \frac{30{,}000}{74} \\
&= 405\,\mathrm{A}
\end{split}
\end{equation}

The per-cell current distributed across 8 parallel strings is:
\begin{equation}
\begin{split}
I_{cell} &= \frac{I_{pack}}{N_p} \\
&= \frac{405}{8} \\
&= 50.6\,\mathrm{A}
\end{split}
\end{equation}

The corresponding discharge C-rate is $50.6 / 5 = 10.1\,\mathrm{C}$.

\begin{table}[H]
\centering
\caption{72\,V Li-ion Battery Pack Electrical Parameters}
\vspace{0.2cm}
\begin{tabular}{@{}lll@{}}
\toprule
\textbf{Parameter} & \textbf{Value} & \textbf{Unit} \\
\midrule
Cell Voltage         & 3.7   & V \\
Configuration        & 20S8P & --- \\
Pack Voltage         & 74    & V \\
Capacity             & 40    & Ah \\
Energy               & 3.0   & kWh \\
Total Cells          & 160   & --- \\
Peak Current         & 405   & A \\
Cell C-Rate          & 10.1  & C \\
\bottomrule
\end{tabular}
\label{tab:battery_electrical}
\end{table}

\paragraph{Mechanical Design.}
The pack enclosure is fabricated from Al 6061 alloy and sealed to IP67 rating for dust and water ingress protection. The 160-cell array is organized into 2 modules of 10S8P each, interconnected via nickel-plated copper busbars. Rubber vibration isolators mount the modules to the chassis to attenuate road-induced harmonics. The estimated pack external dimensions are \SI{400}{\milli\metre} $\times$ \SI{200}{\milli\metre} $\times$ \SI{150}{\milli\metre}, with a total estimated pack mass of approximately \SI{8}{\kilo\gram}.

\paragraph{Thermal Design.}
At a \SI{10}{C} discharge rate, the heat generated per cell is governed by Joule heating:
\begin{equation}
\begin{split}
Q_{cell} &= I_{cell}^2 \cdot R_{int} \\
&= (50.6)^2 \times 0.05 \\
&= 127.9\,\mathrm{W}
\end{split}
\end{equation}

The aggregate peak heat load across all 160 cells is approximately \SI{1350}{\watt} (averaged; peak only during full \SI{30}{\kilo\watt} discharge bursts). To manage this thermal flux, a liquid cold-plate cooling system with aluminium serpentine channels is employed. The coolant is a 50:50 ethylene glycol--water mixture circulated at \SI{3}{\litre\per\minute}. A battery management system (BMS) monitors temperature via NTC thermistors placed at 6 strategic locations across the pack. The permitted operating window is \SIrange{20}{40}{\celsius}, with mandatory shutdown above \SI{60}{\celsius}.

\subsubsection{DC-DC Converter Design}

A bidirectional non-isolated buck-boost converter is selected for the auxiliary DC-DC stage. This topology supports buck mode operation (stepping the \SI{74}{\volt} battery down to a \SI{48}{\volt} auxiliary bus) and boost mode (returning energy from the motor to the battery during regenerative braking). The absence of a transformer reduces overall mass compared to isolated topologies.

\begin{table}[H]
\centering
\caption{DC-DC Converter Specifications}
\vspace{0.2cm}
\begin{tabular}{@{}ll@{}}
\toprule
\textbf{Parameter} & \textbf{Value} \\
\midrule
Input Voltage       & \SI{74}{\volt} \\
Output Voltage      & \SI{48}{\volt} \\
Output Power        & \SI{3000}{\watt} \\
Switching Frequency & \SI{20}{\kilo\hertz} \\
MOSFET Type         & N-channel \\
Topology            & Synchronous Buck-Boost \\
\bottomrule
\end{tabular}
\label{tab:dcdc_specs}
\end{table}

The duty cycle in buck mode is:
\begin{equation}
\begin{split}
D &= \frac{V_{out}}{V_{in}} \\
&= \frac{48}{74} \\
&= 0.649
\end{split}
\end{equation}

The output current at full load is:
\begin{equation}
\begin{split}
I_{out} &= \frac{P_{out}}{V_{out}} \\
&= \frac{3000}{48} \\
&= 62.5\,\mathrm{A}
\end{split}
\end{equation}

Setting the inductor current ripple $\Delta I_L$ to 20\% of the output current:
\begin{equation}
\begin{split}
\Delta I_L &= 0.20 \times 62.5 \\
&= 12.5\,\mathrm{A}
\end{split}
\end{equation}

The required inductance is then:
\begin{equation}
\begin{split}
L &= \frac{V_{out} \cdot (1 - D)}{f_{sw} \cdot \Delta I_L} \\
&= \frac{48 \times 0.351}{20{,}000 \times 12.5} \\
&= \frac{16.85}{250{,}000} \\
&= 67.4\,\mu\mathrm{H} \rightarrow \text{select } L = 68\,\mu\mathrm{H}
\end{split}
\end{equation}

For the output capacitor, limiting the output voltage ripple to 1\% of $V_{out}$ ($\Delta V = \SI{0.48}{\volt}$):
\begin{equation}
\begin{split}
C &= \frac{\Delta I_L}{8 \cdot f_{sw} \cdot \Delta V} \\
&= \frac{12.5}{8 \times 20{,}000 \times 0.48} \\
&= \frac{12.5}{76{,}800} \\
&= 162.8\,\mu\mathrm{F} \rightarrow \text{select } C = 180\,\mu\mathrm{F}
\end{split}
\end{equation}

\subsubsection{Motor Control --- IPM Torque Analysis and FOC}

The Interior Permanent Magnet Synchronous Motor (IPMSM) produces torque through two distinct electromagnetic mechanisms. The total electromagnetic torque is the sum of the magnet alignment torque and the reluctance torque arising from rotor saliency:
\begin{equation}
T_{total} = T_{magnet} + T_{reluctance}
\end{equation}

The magnet torque component is proportional to the quadrature-axis current:
\begin{equation}
\begin{split}
T_{magnet} &= \frac{3}{2} \cdot p \cdot \lambda_f \cdot i_q
\end{split}
\end{equation}

The reluctance torque component arises from the asymmetric inductance along the direct and quadrature axes:
\begin{equation}
\begin{split}
T_{reluctance} &= \frac{3}{2} \cdot p \cdot (L_d - L_q) \cdot i_d \cdot i_q
\end{split}
\end{equation}

For an IPM machine, the quadrature-axis inductance exceeds the direct-axis inductance ($L_q > L_d$), rendering $(L_d - L_q)$ negative. Consequently, the Maximum Torque Per Ampere (MTPA) strategy must set $i_d$ to a negative value to produce positive reluctance torque — a fundamental departure from Surface Permanent Magnet (SPM) machines where $i_d = 0$ is optimal.

The control technique employed is Field Oriented Control (FOC) with Clarke and Park transformations, as detailed in Section~\ref{fig:foc_diagram} of this paper. The FOC algorithm decouples the three-phase AC stator currents into synchronous rotating $d$-axis (flux-producing) and $q$-axis (torque-producing) DC components, enabling independent PI regulation of each axis. The \SI{30}{\kilo\watt} IPM at \SI{72}{\volt} follows the identical FOC architecture documented for the \SI{350}{\kilo\watt} MGU-K in the primary powertrain, scaled to the urban EV operating conditions presented in these reference calculations.
